\documentclass[12pt,a4paper]{article}
\usepackage[utf8]{inputenc}
\usepackage{graphicx}
\usepackage{hyperref}
\usepackage{listings}
\usepackage{xcolor}
\usepackage{geometry}
\usepackage{titlesec}
\usepackage{fancyhdr}
\usepackage{longtable}
\usepackage{enumitem}
\usepackage{setspace}
\usepackage{float} % Required for [H] placement

% Page Geometry
\geometry{a4paper, margin=1in}

% Code snippet styling
\definecolor{codegreen}{rgb}{0,0.6,0}
\definecolor{codegray}{rgb}{0.5,0.5,0.5}
\definecolor{codepurple}{rgb}{0.58,0,0.82}
\definecolor{backcolour}{rgb}{0.95,0.95,0.92}

\lstset{
backgroundcolor=\color{backcolour},   
commentstyle=\color{codegreen},
keywordstyle=\color{magenta},
numberstyle=\tiny\color{codegray},
stringstyle=\color{codepurple},
basicstyle=\ttfamily\footnotesize,
breakatwhitespace=false,         
breaklines=true,                 
captionpos=b,                    
keepspaces=true,                 
numbers=left,                    
numbersep=5pt,                  
showspaces=false,                
showstringspaces=false,
showtabs=false,                  
tabsize=2
}

\begin{document}

% --- Title Page ---
\begin{titlepage}
\centering
\vspace*{3cm}

{\Huge \textbf{Scientific Calculator DevOps Pipeline}} \\
\vspace{2cm}

{\large \textbf{Abstract}} \\
\vspace{0.5cm}
\begin{minipage}{0.8\textwidth}
    \centering
    \textit{This project demonstrates a complete end-to-end CI/CD lifecycle for a Python-based scientific calculator. By integrating Jenkins for orchestration, Docker for containerization, and Ansible for automated deployment, the pipeline ensures that every code change is rigorously tested and seamlessly delivered to a production-ready environment.}
\end{minipage}

\vfill

\begin{flushright}
    \large
    \textbf{Project By:} \\
    \textbf{Name:} Kartikey Dubey \\
    \textbf{Roll Number:} MT2025061
\end{flushright}
\end{titlepage}

\newpage
\tableofcontents
\newpage

% --- GitHub Repository Section ---
\section{Project Metadata}
\textbf{GitHub Repository Link:} \\
\url{https://github.com/kartikeyblaze/SPE-mini-project-calc.git}

\begin{figure}[H]
\centering
\fbox{\includegraphics[width=0.8\textwidth]{images/github-repo.png}}
\caption{Remote Repository Overview}
\end{figure}

% --- Section 1: Introduction ---
\section{Basic Ideas: The What and Why of DevOps}
DevOps is a cultural and technical movement that bridges the historical gap between Software Development (Dev) and IT Operations (Ops). Historically, these two departments functioned as silos, leading to friction during the "hand-over" phase of a project. DevOps solves this by introducing automation, shared responsibility, and continuous feedback loops.

\subsection{Why DevOps?}
\begin{itemize}[noitemsep]
\item \textbf{Speed:} Deliver features and fixes faster through automation.
\item \textbf{Reliability:} Ensure quality via automated unit testing and consistent environments.
\item \textbf{Scalability:} Manage infrastructure as code to handle growth efficiently.
\item \textbf{Collaboration:} Build ownership by involving developers in the deployment lifecycle.
\end{itemize}

% --- Section 2: Tech Stack ---
\section{Tools Used in This Project}
The choice of tools for this project was driven by industry standards and the need for a lightweight, scalable architecture.

\clearpage
\begin{longtable}{|p{3.5cm}|p{3.5cm}|p{7cm}|}
\hline
\textbf{Category} & \textbf{Tool} & \textbf{Purpose} \\
\hline
Source Control & Git / GitHub & Versioning code and hosting the remote repository. \\
\hline
CI/CD Orchestration & Jenkins & Automating the build, test, and deployment stages. \\
\hline
Containerization & Docker & Packaging the application into a portable image. \\
\hline
Artifact Repository & Docker Hub & Storing and sharing built Docker images. \\
\hline
Config Management & Ansible & Automating the deployment on the host. \\
\hline
Programming Language & Python & Logic for the scientific calculator and unit testing. \\
\hline
\end{longtable}

% --- Section 3: Git Phase ---
\section{Phase 1: Source Control with Git}
Source control is the foundation of DevOps. It allows multiple developers to work on the same codebase simultaneously without overwriting each other's work.

\subsection{Implementation \& Commands}
For this project, we followed the \textit{Mainline Development} strategy.
\begin{enumerate}[noitemsep]
\item Initialize the repository: \texttt{git init}
\item Add files: \texttt{git add .}
\item Commit changes: \texttt{git commit -m "Initial commit"}
\item Rename branch: \texttt{git branch -M main}
\item Set remote: \texttt{git remote add origin https://github.com/kartikeyblaze/SPE-mini-project-calc.git}
\item Push to GitHub: \texttt{git push -u origin main}
\end{enumerate}

\begin{figure}[H]
\centering
\includegraphics[width=0.8\textwidth]{images/git-push.png}
\caption{Successful Code Push to GitHub}
\end{figure}

% --- Section 4: Jenkins Phase ---
\section{Phase 2: Jenkins Orchestration}
Jenkins serves as the "brain" of the operation, triggering builds automatically upon every code change detected via webhooks.

\subsection{Setup \& Configuration}
\begin{enumerate}[noitemsep]
\item \textbf{Expose Jenkins with ngrok:} Run \texttt{ngrok http 8080} to get a public URL.
\begin{figure}[H]
  \centering
  \includegraphics[width=0.8\textwidth]{images/ngrok.png}
  \caption{Exposing local Jenkins via ngrok}
\end{figure}

\item \textbf{GitHub Webhook:} Added the ngrok URL followed by \texttt{/github-webhook/} in GitHub Settings. 
\begin{figure}[H]
\centering
    \includegraphics[width=0.8\textwidth]{images/github-webhook.png}
    \caption{Connectivity and Webhook Setup}
\end{figure}

\item \textbf{Docker Hub Credentials:} Stored as "Username with password" with ID \texttt{DockerHubCred}.
\begin{figure}[H]
\centering
    \includegraphics[width=0.8\textwidth]{images/jenkins-docker-cred.png}
    \caption{Docker Hub Credentials configuration}
\end{figure}
\item \textbf{Mail Notification:} Configured via SMTP (port 587/465) for status alerts.
\begin{figure}[H]
\centering
    \includegraphics[width=0.8\textwidth]{images/jenkins-smtp.png}
    \caption{SMTP Server settings for Notifications}
\end{figure}
\end{enumerate}

% --- Section 5: Jenkinsfile ---
\section{Phase 3: Pipeline as Code (Jenkinsfile)}
A key practice is "Pipeline as Code." The \texttt{Jenkinsfile} ensures the build process is repeatable and version-controlled.

\begin{lstlisting}[language=Groovy, caption=Scientific Calculator Jenkinsfile]
pipeline {
agent any
environment {
    DOCKER_HUB_USER = 'kartikeyblaze'
    IMAGE_NAME = 'scientific-calculator'
}
stages {
    stage('Checkout') { steps { checkout scm } }
    stage('Unit Testing') { steps { sh 'python3 -m unittest test_calculator.py' } }
    stage('Build Docker Image') { 
        steps {
            sh "docker build -t ${DOCKER_HUB_USER}/${IMAGE_NAME}:${env.BUILD_ID} ."
        }
    }
    stage('Push to Docker Hub') {
        steps {
            script {
                withCredentials([usernamePassword(credentialsId: 'DockerHubCred', ...)]) {
                    sh 'docker push ...'
                }
            }
        }
    }
    stage('Run Ansible Playbook') { steps { ansiblePlaybook(playbook: 'deploy.yml', inventory: 'inventory') } }
}
post { always { script { mail to: EMAIL, subject: "Build Result", body: "Build Status: ${currentBuild.result}" } } }
}
\end{lstlisting}

\section*{Block-wise Explanation}
\begin{itemize}[leftmargin=*]
\item \textbf{pipeline \{ agent any \}:} Directs Jenkins to run this pipeline on any available agent.
\item \textbf{environment \{ \ldots \}:} Defines global variables like \texttt{DOCKER\_HUB\_USER} and \texttt{IMAGE\_NAME} used throughout the stages.
\item \textbf{triggers \{ githubPush() \}:} Tells Jenkins to start the build whenever a push event is received from the GitHub webhook.
\item \textbf{stage('Unit Testing'):} Runs \texttt{unittest} on \texttt{test\_calculator.py}. If tests fail, the pipeline stops.
\item \textbf{stage('Build Docker Image'):} Builds the image using the local \texttt{Dockerfile} and tags it with the unique build ID.
\item \textbf{stage('Push to Docker Hub'):} Securely logs into Docker Hub and pushes the images.
\item \textbf{stage('Run Ansible Playbook'):} Triggers the automated deployment phase.
\item \textbf{post \{ always \{ \ldots \} \}:} Sends an email notification regardless of whether the build succeeded or failed.
\end{itemize}

\begin{figure}[H]
\centering
\includegraphics[width=0.8\textwidth]{images/jenkins-pipeline-config.png}
\caption{Jenkins Dashboard showing the created Pipeline Job}
\end{figure}

\begin{figure}[H]
\centering
\includegraphics[width=0.8\textwidth]{images/jenkins-pipeline.png}
\caption{Jenkins Pipeline Stage View }
\end{figure}

% --- Section 6: Docker ---
\section{Phase 4: Containerization with Docker}
Docker solves the "Dependency Hell" problem by packaging the Python runtime and script into a secure layer.

\subsection{Dockerfile Breakdown}
\begin{lstlisting}[language=bash]
FROM python:3.9-slim  # Lightweight base
WORKDIR /app          # Working directory
COPY Calculator.py .  # App logic
CMD ["python", "Calculator.py"] # Run instruction
\end{lstlisting}

\subsection{Detailed Working of the Dockerfile}
The \texttt{Dockerfile} implementation uses a layered architecture to ensure a small footprint and fast builds.
\begin{itemize}
    \item \textbf{Base Image (\texttt{FROM}):} We use \texttt{python:3.9-slim}. The "slim" tag indicates a minimal image that contains only the essential packages, reducing image size and vulnerabilities.
    \item \textbf{Working Directory (\texttt{WORKDIR}):} Creates \texttt{/app} inside the container and sets it as the active directory for subsequent commands.
    \item \textbf{Copying Source Code (\texttt{COPY}):} Copies the \texttt{Calculator.py} script from the local Jenkins workspace into the container's current working directory.
    \item \textbf{Execution Command (\texttt{CMD}):} Specifies the default command to execute when the container starts, invoking the Python interpreter to run the script.
\end{itemize}

\begin{figure}[H]
\centering
\includegraphics[width=0.8\textwidth]{images/dockerhub-uploaded.png}
\caption{Artifact Storage in Docker Hub}
\end{figure}

% --- Section 7: Ansible ---
\section{Phase 5: Deployment with Ansible}
Ansible ensures the host machine is in the desired state by pulling the latest image and managing the container life cycle.

\subsection{Detailed Playbook Breakdown}
The \texttt{deploy.yml} file is an Ansible playbook that defines the desired state of the deployment server. 
\begin{lstlisting}[language=bash, caption=Ansible Deployment Playbook (deploy.yml)]
---
- name: Pull Docker Image and Start Container
  hosts: localhost
  connection: local
  tasks:
    - name: Pull Docker Image
      docker_image:
        name: "kartikeyblaze/scientific-calculator:latest"
        source: pull
    - name: Remove existing container
      docker_container:
        name: "calculator"
        state: absent
    - name: Running container
      shell: docker run -it -d --name calculator kartikeyblaze/scientific-calculator:latest
\end{lstlisting}

\begin{itemize}
    \item \textbf{Targeting (\texttt{hosts} \& \texttt{connection}):} Since we are deploying to the local machine, we set \texttt{hosts: localhost} and \texttt{connection: local} to bypass SSH.
    \item \textbf{Image Retrieval (\texttt{docker\_image}):} Ensures that the latest version of the image is pulled from Docker Hub.
    \item \textbf{Clean-up (\texttt{docker\_container}):} Explicitly removes any existing \texttt{calculator} container to prevent naming conflicts.
    \item \textbf{Execution (\texttt{shell}):} Starts the container in detached mode (\texttt{-d}), allowing it to persist in the background.
\end{itemize}

\begin{figure}[H]
\centering
\includegraphics[width=0.8\textwidth]{images/docker-ps.png}
\caption{Successful Final Deployment Status}
\end{figure}

% --- Section 8: Email ---
\section{Email Notifications}
Automated emails are sent at the end of every build to keep the team informed of the production health.

\begin{figure}[H]
\centering
\includegraphics[width=0.8\textwidth]{images/email-failure.png}
\vspace{0.5cm}
\includegraphics[width=0.8\textwidth]{images/email-success.png}
\caption{Email notification results (Failure and Success)}
\end{figure}

% --- Section 9: Conclusion ---
\section{Conclusion}
This project successfully establishes a robust, automated pipeline for a scientific calculator. By eliminating manual steps, we have significantly reduced the risk of deployment errors.

\end{document}
